% Generated by scripts/build_prompt_appendix.py. Do not edit by hand.
% Every block is the literal prompt text, verified against the stored
% per-record prompt hashes; see review/audits/PROMPT_FIDELITY_AUDIT.json.

\paragraph{Reading these blocks.}
Placeholders in braces are filled per record. Every other character is
literal, including the runs of three dots inside the output template: they
are part of the prompt the model received and are not an omission by us.
At Stage~2 the image is not a token in the prompt text. It is passed as a
separate image item in the model's chat template, and the processor inserts
its own image tokens. The Stage~1 text opens with the literal string
\texttt{<image>}.

\paragraph{Stage 1.}
The image accompanies the following text.

\begin{quote}\small\ttfamily\frenchspacing\setlength{\parskip}{0pt}
\textless image\textgreater \\
Question: \{question\}\\
Choices:\\
\hspace*{1ex}\hspace*{1ex}(A) \{c0\}\\
\hspace*{1ex}\hspace*{1ex}(B) \{c1\}\\
\hspace*{1ex}\hspace*{1ex}(C) \{c2\}\\
\hspace*{1ex}\hspace*{1ex}(D) \{c3\}\\
Instruction: Select the best option and answer with a single letter (A/B/C/D).
\end{quote}

Description arms append one blank line and then this field, which closes
the Stage~1 prompt; the other arms end after the instruction line.

\begin{quote}\small\ttfamily\frenchspacing\setlength{\parskip}{0pt}
Caption: \{caption\}
\end{quote}

The answer is taken from the classifier head, so the letter instruction
frames the task without being parsed (Section~\ref{sec:setup-model}).

\paragraph{Stage 2 with the gold answer.}
This is the prompt behind every primary RQ2 generation and behind the
wrong-answer control, which differs only in the answer index and text
it supplies.

\begin{quote}\small\ttfamily\frenchspacing\setlength{\parskip}{0pt}
You are an expert visual commonsense reasoner.\\
\mbox{}\\
You see an image and a multiple-choice question with four options.\\
Your task is to explain, step by step, why the given answer is correct,\\
using only what can reasonably be inferred from the image and question.\\
\mbox{}\\
Question: \{question\}\\
\mbox{}\\
Choices:\\
\hspace*{1ex}\hspace*{1ex}(0) \{c0\}\\
\hspace*{1ex}\hspace*{1ex}(1) \{c1\}\\
\hspace*{1ex}\hspace*{1ex}(2) \{c2\}\\
\hspace*{1ex}\hspace*{1ex}(3) \{c3\}\\
\mbox{}\\
The correct answer is choice index \{answer\_idx\}: \textquotedbl \{answer\}\textquotedbl .\\
\mbox{}\\
Explain your reasoning step by step, focusing only on what can be seen\\
in the image (and what follows directly from it).\\
Then give a short final justification.\\
\mbox{}\\
Use the following format exactly:\\
\mbox{}\\
\textless reasoning\textgreater \\
...detailed step-by-step reasoning about the image, question, and answer...\\
\textless /reasoning\textgreater \\
\textless final\textgreater \\
...short 2-3 sentence justification...\\
\textless /final\textgreater 
\end{quote}

The two description arms append one blank line and then this field, so the
description is the last thing in the prompt.

\begin{quote}\small\ttfamily\frenchspacing\setlength{\parskip}{0pt}
Image caption: \{caption\}
\end{quote}

\paragraph{Training and generation prompts.}
Stage~2 was trained with a template close to the one above but not identical
to it. All four arms used the same training template. It asks why \emph{a}
given answer is correct rather than \emph{the} given answer, and it adds one
instruction, \enquote{Continue until the explanation is complete and
well-formed}, that the generation prompt does not carry. In the two
description arms the training prompt also appended a different description
block from the one used at generation: it labelled the description as context
and told the model not to reuse it.

\begin{quote}\small\ttfamily\frenchspacing\setlength{\parskip}{0pt}
Caption (context only — do NOT copy this text in your reasoning): \{caption\}\\
Your explanation MUST NOT reuse the caption. Use only the image + question + answer.
\end{quote}

Both templates request a two- to three-sentence final justification, whereas
the training target places the last sentence of the human rationale in the
\texttt{<final>} field (Section~\ref{sec:setup-prompts}). The analysis
therefore evaluates the non-empty final span the model actually generated,
whatever its sentence count.

\paragraph{Stage 2 with the predicted answer.}
The secondary predicted-answer regime used this prompt. It frames the task
as reflection on the model’s own decision, supplies the Stage~1 answer
rather than the gold answer, and adds the instruction to continue until the
explanation is complete. It is not the primary generation prompt.

\begin{quote}\small\ttfamily\frenchspacing\setlength{\parskip}{0pt}
You are reflecting on your own decision in a visual question answering task.\\
\mbox{}\\
You see an image and a multiple-choice question with four options.\\
Your task is to explain, step by step, why the answer you chose is (or seems)\\
correct, using only what can reasonably be inferred from the image and question.\\
\mbox{}\\
Question: \{question\}\\
\mbox{}\\
Choices:\\
\hspace*{1ex}\hspace*{1ex}(0) \{c0\}\\
\hspace*{1ex}\hspace*{1ex}(1) \{c1\}\\
\hspace*{1ex}\hspace*{1ex}(2) \{c2\}\\
\hspace*{1ex}\hspace*{1ex}(3) \{c3\}\\
\mbox{}\\
You picked choice index \{answer\_idx\}: \textquotedbl \{answer\}\textquotedbl .\\
\mbox{}\\
Explain your reasoning step by step, focusing only on what can be seen\\
in the image (and what follows directly from it).\\
Then give a short final justification.\\
Continue until the explanation is complete and well-formed.\\
\mbox{}\\
Use the following format exactly:\\
\mbox{}\\
\textless reasoning\textgreater \\
...detailed step-by-step reasoning about the image, question, and answer...\\
\textless /reasoning\textgreater \\
\textless final\textgreater \\
...short 2–3 sentence justification...\\
\textless /final\textgreater 
\end{quote}

The two description arms append this training block. The primary generation
prompt uses a different description field.
The other two arms receive no description block at all.

\begin{quote}\small\ttfamily\frenchspacing\setlength{\parskip}{0pt}
Caption (context only — do NOT copy this text in your reasoning): \{caption\}\\
Your explanation MUST NOT reuse the caption. Use only the image + question + answer.
\end{quote}

\paragraph{Cross-Modal Consistency judge.}
The judge receives this text and nothing else. There is no reasoning or
final scaffold: the \texttt{Explanation} field holds the eligible final
rationale span on its own, so the generated reasoning span never reaches the
judge, and neither does the image, the arm, or the answer supplied to
Stage~2. The judge emits one token, constrained to the four literal digits,
so every record returns a choice index and no parse can fail.

\begin{quote}\small\ttfamily\frenchspacing\setlength{\parskip}{0pt}
You are judging a multiple-choice question without access to its image.\\
Another system supplied an explanation. Use only the question, choices, and explanation.\\
\mbox{}\\
Question: \{question\}\\
\mbox{}\\
Choices:\\
(0) \{c0\}\\
(1) \{c1\}\\
(2) \{c2\}\\
(3) \{c3\}\\
\mbox{}\\
Explanation:\\
\{rationale\_final\}\\
\mbox{}\\
Which choice is best supported? Answer with exactly one digit: 0, 1, 2, or 3.\\
Answer:
\end{quote}

\paragraph{Description generation.}
Descriptions were produced from the \mbox{polygon-marked} frames through the
OpenAI batch interface with the \texttt{gpt-4o} alias, temperature 0.2, and a
limit of 120 output tokens. The image was supplied inline as a base64 JPEG.
No dated model snapshot was pinned, which
Section~\ref{sec:discussion-limitations} records as a reproducibility
limitation.

\begin{quote}\small\ttfamily\frenchspacing\setlength{\parskip}{0pt}
System: You are a vision-language model tasked with generating factual, concise scene descriptions based on labeled image elements such as person\_0 or object\_1. Avoid speculation and emotional language.
\end{quote}

\begin{quote}\small\ttfamily\frenchspacing\setlength{\parskip}{0pt}
User: Describe this scene based on labeled elements (e.g., person\_0, object\_1). Focus on positions, actions, and interactions without interpretation. Be concise and objective.
\end{quote}
